% ABSTRACT PAGE
\thispagestyle{empty}
\newpage
\cleardoublepage\phantomsection
\addcontentsline{toc}{chapter}{Abstract}\index{Abstract}
\begin{center}
{\Large \bf Abstract}\\
\end{center}
\vspace{10pt}

\noindent
In modern software engineering and computer science pedagogy, source code quality, maintainability, and readability are as critical as algorithmic correctness and runtime performance. However, conventional industrial compilers treat code comments as trivial whitespace tokens and immediately discard them during the lexical scanning phase. Consequently, software systems frequently suffer from severe documentation debt, making complex logic difficult to comprehend, verify, and maintain across collaborative and academic environments.

This project presents the design, theoretical formulation, and end-to-end implementation of \textbf{LightAngo}, an ahead-of-time (AOT) compiler for a custom strongly-typed high-level imperative programming language. Developed in modern C++ (C++20 standard) under the Linux environment, LightAngo realizes a full six-phase compilation pipeline: an enhanced state-machine lexical scanner, a recursive descent abstract syntax tree (AST) parser, a scoped semantic analyzer with symbol table resolution, an intermediate representation synthesizer, machine-independent optimization passes, and a backend code generator emitting native 64-bit x86 (x86-64) NASM assembly directly linkable into executable ELF binaries.

A defining innovation of LightAngo is its integrated \textit{Grammar-Level Documentation Enforcement Mechanism}. Rather than discarding source comments as irrelevant artifacts, the lexical and syntactic parsers track preceding comment tokens attached to structural grammar productions. The compiler enforces strict semantic policies requiring clear, meaningful documentation headers before every loop construct and function declaration. If mandatory comment specifications are absent or malformed, the compiler halts with precise diagnostic errors pointing to line and column locations, preventing undocumented logic from reaching production binaries.

Additionally, the compiler features comprehensive semantic type verification, variable lifetime checks, control-flow jump resolution, and expression stack flattening. Comprehensive experimental evaluations demonstrate that LightAngo successfully generates accurate, optimized assembly for complex arithmetic, conditional branching, nested iteration, and subroutines, while demonstrating that documentation-aware compilation enforces disciplined programming habits without imposing measurable runtime performance overhead.

\vspace{1.0cm}
\noindent
\textbf{Keywords:} Compiler Design, Lexical Analysis, Abstract Syntax Tree, Semantic Analysis, x86-64 Code Generation, NASM, Mandatory Documentation, Code Readability, Software Quality, Static Analysis.
\newpage
\cleardoublepage
